\documentclass[11pt]{article}
\newcommand{\ListaTitulo}{Gabarito 03: Dispersão: AT, DP, variância e CV}
% Preâmbulo compartilhado: MATF14, Listas de Exercícios 2026
% Usado por ListaNN.tex e GabaritoNN.tex via % Preâmbulo compartilhado: MATF14, Listas de Exercícios 2026
% Usado por ListaNN.tex e GabaritoNN.tex via % Preâmbulo compartilhado: MATF14, Listas de Exercícios 2026
% Usado por ListaNN.tex e GabaritoNN.tex via \input{preamble}

\usepackage[utf8]{inputenc}
\usepackage[T1]{fontenc}
\usepackage[brazil]{babel}
\usepackage{fullpage}
\usepackage{amsmath,amssymb}
\usepackage{enumitem}
\usepackage{xcolor}
\usepackage{fancyhdr}
\usepackage{titlesec}
\usepackage{listings}
\usepackage{hyperref}
\usepackage{booktabs}
\usepackage{graphicx}
\usepackage{tikz}

\definecolor{matf14azul}{HTML}{0D3B54}
\definecolor{matf14dourado}{HTML}{C9971E}

\titleformat{\section}{\normalfont\Large\bfseries\color{matf14azul}}{\thesection}{1em}{}
\titleformat{\subsection}{\normalfont\large\bfseries\color{matf14azul}}{\thesubsection}{1em}{}

\providecommand{\ListaTitulo}{Lista de Exercícios}

\setlength{\headheight}{14pt}
\pagestyle{fancy}
\fancyhf{}
\lhead{\small MATF14: Estatística Econômica I}
\rhead{\small \ListaTitulo}
\cfoot{\thepage}
\renewcommand{\headrulewidth}{0.4pt}

\lstset{
  basicstyle=\ttfamily\small,
  breaklines=true,
  frame=single,
  columns=fullflexible,
  backgroundcolor=\color{gray!8},
  commentstyle=\color{matf14dourado},
  keywordstyle=\color{matf14azul}\bfseries,
  language=R
}

\newcounter{questaocounter}
\newenvironment{questao}[1]{%
  \stepcounter{questaocounter}%
  \bigskip\noindent\textbf{Questão #1.}\ %
}{\par}

\newcommand{\cabecalho}[2]{%
  \begin{center}
    {\Large\bfseries\color{matf14azul} #1}\\[0.3em]
    {\large #2}
  \end{center}
  \vspace{0.5em}
  \noindent\rule{\linewidth}{0.6pt}
  \vspace{0.5em}
}

\newenvironment{solucao}{%
  \par\smallskip\noindent\textit{\color{matf14dourado!80!black}Solução.}\ %
}{\hfill$\blacksquare$\par}


\usepackage[utf8]{inputenc}
\usepackage[T1]{fontenc}
\usepackage[brazil]{babel}
\usepackage{fullpage}
\usepackage{amsmath,amssymb}
\usepackage{enumitem}
\usepackage{xcolor}
\usepackage{fancyhdr}
\usepackage{titlesec}
\usepackage{listings}
\usepackage{hyperref}
\usepackage{booktabs}
\usepackage{graphicx}
\usepackage{tikz}

\definecolor{matf14azul}{HTML}{0D3B54}
\definecolor{matf14dourado}{HTML}{C9971E}

\titleformat{\section}{\normalfont\Large\bfseries\color{matf14azul}}{\thesection}{1em}{}
\titleformat{\subsection}{\normalfont\large\bfseries\color{matf14azul}}{\thesubsection}{1em}{}

\providecommand{\ListaTitulo}{Lista de Exercícios}

\setlength{\headheight}{14pt}
\pagestyle{fancy}
\fancyhf{}
\lhead{\small MATF14: Estatística Econômica I}
\rhead{\small \ListaTitulo}
\cfoot{\thepage}
\renewcommand{\headrulewidth}{0.4pt}

\lstset{
  basicstyle=\ttfamily\small,
  breaklines=true,
  frame=single,
  columns=fullflexible,
  backgroundcolor=\color{gray!8},
  commentstyle=\color{matf14dourado},
  keywordstyle=\color{matf14azul}\bfseries,
  language=R
}

\newcounter{questaocounter}
\newenvironment{questao}[1]{%
  \stepcounter{questaocounter}%
  \bigskip\noindent\textbf{Questão #1.}\ %
}{\par}

\newcommand{\cabecalho}[2]{%
  \begin{center}
    {\Large\bfseries\color{matf14azul} #1}\\[0.3em]
    {\large #2}
  \end{center}
  \vspace{0.5em}
  \noindent\rule{\linewidth}{0.6pt}
  \vspace{0.5em}
}

\newenvironment{solucao}{%
  \par\smallskip\noindent\textit{\color{matf14dourado!80!black}Solução.}\ %
}{\hfill$\blacksquare$\par}


\usepackage[utf8]{inputenc}
\usepackage[T1]{fontenc}
\usepackage[brazil]{babel}
\usepackage{fullpage}
\usepackage{amsmath,amssymb}
\usepackage{enumitem}
\usepackage{xcolor}
\usepackage{fancyhdr}
\usepackage{titlesec}
\usepackage{listings}
\usepackage{hyperref}
\usepackage{booktabs}
\usepackage{graphicx}
\usepackage{tikz}

\definecolor{matf14azul}{HTML}{0D3B54}
\definecolor{matf14dourado}{HTML}{C9971E}

\titleformat{\section}{\normalfont\Large\bfseries\color{matf14azul}}{\thesection}{1em}{}
\titleformat{\subsection}{\normalfont\large\bfseries\color{matf14azul}}{\thesubsection}{1em}{}

\providecommand{\ListaTitulo}{Lista de Exercícios}

\setlength{\headheight}{14pt}
\pagestyle{fancy}
\fancyhf{}
\lhead{\small MATF14: Estatística Econômica I}
\rhead{\small \ListaTitulo}
\cfoot{\thepage}
\renewcommand{\headrulewidth}{0.4pt}

\lstset{
  basicstyle=\ttfamily\small,
  breaklines=true,
  frame=single,
  columns=fullflexible,
  backgroundcolor=\color{gray!8},
  commentstyle=\color{matf14dourado},
  keywordstyle=\color{matf14azul}\bfseries,
  language=R
}

\newcounter{questaocounter}
\newenvironment{questao}[1]{%
  \stepcounter{questaocounter}%
  \bigskip\noindent\textbf{Questão #1.}\ %
}{\par}

\newcommand{\cabecalho}[2]{%
  \begin{center}
    {\Large\bfseries\color{matf14azul} #1}\\[0.3em]
    {\large #2}
  \end{center}
  \vspace{0.5em}
  \noindent\rule{\linewidth}{0.6pt}
  \vspace{0.5em}
}

\newenvironment{solucao}{%
  \par\smallskip\noindent\textit{\color{matf14dourado!80!black}Solução.}\ %
}{\hfill$\blacksquare$\par}


\begin{document}

\cabecalho{Gabarito 03}{Amplitude total, desvio padrão, variância e coeficiente de variação (Cap. 1, Aula 06)}

\begin{questao}{1} \textbf{(Cálculo direto: retorno de dois fundos de investimento)}
\begin{solucao}
\begin{enumerate}[label=(\alph*)]
  \item Fundo A: $\bar x_A = (1{,}0+1{,}2+0{,}9+1{,}1+0{,}8)/5 = 5{,}0/5 = 1{,}0$.
    Fundo B: $\bar x_B = (-2{,}0+4{,}0+0{,}5+3{,}0-0{,}5)/5 = 5{,}0/5 = 1{,}0$. Médias idênticas.
  \item Fundo A: desvios $0, 0{,}2, -0{,}1, 0{,}1, -0{,}2$; quadrados $0, 0{,}04, 0{,}01, 0{,}01,
    0{,}04$; soma $=0{,}10$; $\mathrm{Var}_A = 0{,}10/5=0{,}02$; $dp_A=\sqrt{0{,}02}\approx
    0{,}141$.
    Fundo B: desvios $-3, 3, -0{,}5, 2, -1{,}5$; quadrados $9, 9, 0{,}25, 4, 2{,}25$; soma
    $=24{,}5$; $\mathrm{Var}_B=24{,}5/5=4{,}9$; $dp_B=\sqrt{4{,}9}\approx 2{,}214$.
  \item Um investidor avesso a risco prefere o Fundo A: mesma média de retorno, mas desvio padrão
    muito menor ($0{,}141$ contra $2{,}214$), menos incerteza sobre o retorno mensal exato.
\end{enumerate}
\end{solucao}
\end{questao}

\begin{questao}{2} \textbf{(Coeficiente de variação: comparando variáveis em unidades diferentes)}
\begin{solucao}
\begin{enumerate}[label=(\alph*)]
  \item Não é correto comparar diretamente os desvios padrão: eles estão em unidades diferentes
    ("número de funcionários" vs. "milhões de reais"), com médias em escalas muito diferentes
    (45 vs. 12). Um desvio padrão de 18 pessoas em torno de uma média de 45 é proporcionalmente
    muito maior que um desvio padrão de 3 milhões em torno de uma média de 12 milhões, mas isso
    só se vê normalizando pela média, não comparando os desvios padrão brutos.
  \item $CV_{\text{func}} = 18/45 = 0{,}40$ (40\%). $CV_{\text{fat}} = 3/12 = 0{,}25$ (25\%). O
    número de funcionários é, de fato, \textbf{relativamente} mais disperso (CV maior) que o
    faturamento, mas essa é uma conclusão diferente, e mais confiável, do que comparar os desvios
    padrão brutos em (a).
\end{enumerate}
\end{solucao}
\end{questao}

\begin{questao}{3} \textbf{(Amplitude total vs. desvio padrão: sensibilidade a outliers)}
\begin{solucao}
\begin{enumerate}[label=(\alph*)]
  \item Time 1: $AT_1 = 23-19=4$. Time 2: $AT_2 = 90-19=71$.
  \item Time 1: $\bar x_1 = (20+22+21+23+19+20+21)/7 = 146/7\approx 20{,}86$; calculando,
    $dp_1\approx 1{,}25$. Time 2 (com o 90): $\bar x_2 = (20+21+20+22+19+21+90)/7 = 213/7\approx
    30{,}43$; $dp_2\approx 24{,}9$. O desvio padrão do Time 2 dispara, assim como $AT$, mas de
    forma ainda mais proporcionalmente distorcida, pois o quadrado do desvio amplia o efeito de um
    valor muito afastado. Tanto $AT$ quanto $dp$ (calculado sobre \emph{todos} os pontos) são
    sensíveis a um único valor atípico; nenhuma das duas é robusta.
  \item O $IQR$ (Aula 05), por depender só dos quartis centrais e ignorar os valores extremos
    (inclusive o 90, que ficaria fora da "caixa" central), é a medida de dispersão recomendada
    aqui, ela descreveria a consistência das vendas "normais" do Time 2 sem ser distorcida pelo
    valor suspeito, que pode ser investigado separadamente.
\end{enumerate}
\end{solucao}
\end{questao}

\begin{questao}{4} \textbf{(Dedução: a variância é sempre não negativa)}
\begin{solucao}
\begin{enumerate}[label=(\alph*)]
  \item Cada termo $(x_i - \bar x)^2$ é o quadrado de um número real, logo $(x_i-\bar x)^2 \ge 0$
    para todo $i$. A soma $\sum_{i=1}^n (x_i-\bar x)^2$ é, portanto, uma soma de termos não
    negativos, logo é ela mesma não negativa. Dividir por $n>0$ preserva o sinal, logo
    $\mathrm{Var}(X) = \frac{1}{n}\sum_i (x_i-\bar x)^2 \ge 0$.
  \item ($\Leftarrow$) Se $x_1=\cdots=x_n=c$ para alguma constante $c$, então $\bar x = c$ e cada
    $(x_i-\bar x)^2 = (c-c)^2=0$, logo $\mathrm{Var}(X)=0$.
    ($\Rightarrow$) Se $\mathrm{Var}(X)=0$, então $\sum_{i=1}^n (x_i-\bar x)^2 = 0$. Como essa é
    uma soma de $n$ termos não negativos (item a) que soma zero, cada termo individual deve ser
    zero (uma soma de números $\ge 0$ só é zero se todas as parcelas forem zero), logo
    $(x_i-\bar x)^2=0$ para todo $i$, ou seja, $x_i=\bar x$ para todo $i$: todos os valores são
    iguais entre si (e iguais à própria média).
\end{enumerate}
\end{solucao}
\end{questao}

\begin{questao}{5} \textbf{(Dedução: fórmula alternativa da variância)}
\begin{solucao}
Expandindo o quadrado dentro do somatório:
$$
\sum_{i=1}^n (x_i-\bar x)^2 = \sum_{i=1}^n \left(x_i^2 - 2x_i\bar x + \bar x^2\right)
= \sum_{i=1}^n x_i^2 - 2\bar x\sum_{i=1}^n x_i + n\bar x^2.
$$
Como $\bar x$ é constante (não depende de $i$), ela sai dos somatórios sem alteração. Usando
$\sum_{i=1}^n x_i = n\bar x$ (definição de média):
$$
\sum_{i=1}^n (x_i-\bar x)^2 = \sum_{i=1}^n x_i^2 - 2\bar x (n\bar x) + n\bar x^2
= \sum_{i=1}^n x_i^2 - 2n\bar x^2 + n\bar x^2 = \sum_{i=1}^n x_i^2 - n\bar x^2.
$$
Dividindo por $n$ ambos os lados:
$$
\mathrm{Var}(X) = \frac{1}{n}\sum_{i=1}^n (x_i-\bar x)^2 = \frac{1}{n}\sum_{i=1}^n x_i^2 - \bar x^2,
$$
como queríamos mostrar.

\textbf{Conferindo com o Fundo A} (Questão 1): $x_i^2 = 1{,}00, 1{,}44, 0{,}81, 1{,}21, 0{,}64$,
soma $=5{,}10$, logo $\frac{1}{5}(5{,}10) - (1{,}0)^2 = 1{,}02 - 1{,}00 = 0{,}02$, exatamente o
mesmo valor de $\mathrm{Var}_A$ obtido na Questão 1(b) pela fórmula original.
\end{solucao}
\end{questao}

\end{document}
